\clearpage\pagestyle{otMain}\setlength{\otBlocEspacement}{2.5mm}\fontsize{11}{15}\selectfont\corrigesChapitre{18}{Organiser et interpréter des données}
\begin{blocLivre}{cor-18-1}\begin{corrige}{1}{Un effectif total}\demarche{Additionner les catégories.}
$7+9+4=\res{20}$ réponses. Chaque réponse est comptée une seule fois.
\end{corrige}\end{blocLivre}
\begin{blocLivre}{cor-18-2}\begin{corrige}{2}{Une fréquence}\demarche{Diviser l’effectif de la catégorie par le total.}
$10/25=2/5=\res{0{,}4=40\,\%}$. Le total de référence est 25.
\end{corrige}\end{blocLivre}
\begin{blocLivre}{cor-18-3}\begin{corrige}{3}{Une moyenne directe}\demarche{Diviser la somme par le nombre de données.}
La somme vaut 44. La moyenne est $44:4=\res{11}$.
\end{corrige}\end{blocLivre}
\begin{blocLivre}{cor-18-4}\begin{corrige}{4}{Lire des barres}\demarche{Lire le rôle de chaque axe.}
La barre d’abscisse 2 atteint l’effectif $\res4$. L’axe horizontal donne les livres et le vertical les élèves.
\label{fig-cor-18-4}\figureGeo[corrige]{stat-4}
\end{corrige}\end{blocLivre}
\begin{blocLivre}{cor-18-5}\begin{corrige}{5}{Choisir un graphique}\demarche{Relier la représentation à l’objectif.}
Pour les sports, un $\res{\text{diagramme en barres}}$ permet de comparer les catégories. Pour la température selon l’heure, un $\res{\text{graphique cartésien}}$ montre l’évolution.
\end{corrige}\end{blocLivre}
\begin{blocLivre}{cor-18-6}\begin{corrige}{6}{Organiser une liste}\demarche{Compter puis contrôler le total.}
Les effectifs sont $3,4,1$, total 8. Les fréquences sont $\res{3/8,\ 1/2,\ 1/8}$, soit 37,5 \%, 50 \% et 12,5 \%. Leur somme vaut 100 \%.
\end{corrige}\end{blocLivre}
\begin{blocLivre}{cor-18-7}\begin{corrige}{7}{Construire des secteurs}\demarche{Traduire les parts du total en parts du tour.}
Les fréquences sont 50 \%, 30 \% et 20 \%. En multipliant par $360^\circ$, on obtient $\res{180^\circ,\ 108^\circ,\ 72^\circ}$. Le total est 360 degrés.
\label{fig-cor-18-7}\figureGeo[corrige]{stat-7}
\end{corrige}\end{blocLivre}
\begin{blocLivre}{cor-18-8}\begin{corrige}{8}{Lire un secteur}\demarche{Rapporter l’angle au tour complet.}
$90/360=1/4$. Un quart de 48 vaut $\res{12}$ personnes.
\label{fig-cor-18-8}\figureGeo[corrige]{stat-8}
\end{corrige}\end{blocLivre}
\begin{blocLivre}{cor-18-9}\begin{corrige}{9}{Une moyenne décimale}\demarche{Additionner dans la même unité.}
La somme vaut 12 cm. La moyenne est $12:4=\res{3\ \mathrm{cm}}$, comprise entre la plus petite et la plus grande mesure. Avec 0,5 cm de plus pour chacune, la somme augmente de $4\times0{,}5=2$ cm : la nouvelle moyenne vaut $14:4=\res{3{,}5\ \mathrm{cm}}$. Elle augmente aussi de 0,5 cm.
\end{corrige}\end{blocLivre}
\begin{blocLivre}{cor-18-10}\begin{corrige}{10}{Une donnée extrême}\demarche{Comparer les sommes à effectif constant.}
La première moyenne vaut $\res5$. La seconde vaut $69:3=\res{23}$. La valeur 60 tire fortement la moyenne vers le haut.
\end{corrige}\end{blocLivre}
\begin{blocLivre}{cor-18-11}\begin{corrige}{11}{Retrouver une donnée}\demarche{Retrouver d’abord la somme totale.}
La somme des cinq doit valoir $5\times12=60$. La somme connue vaut 47. Le dernier nombre est $60-47=\res{13}$.
\end{corrige}\end{blocLivre}
\begin{blocLivre}{cor-18-12}\begin{corrige}{12}{Comparer des groupes}\demarche{Distinguer effectif et fréquence.}
B a plus de réussites : 18 contre 12. Mais les taux sont $12/20=60\,\%$ et $18/40=45\,\%$. Le meilleur taux est celui de $\res A$.
\end{corrige}\end{blocLivre}
\begin{blocLivre}{cor-18-13}\begin{corrige}{13}{Même moyenne, profils différents}\demarche{Revenir aussi aux valeurs individuelles.}
Les deux sommes valent 20, donc les deux moyennes valent $\res5$. La première série est constante, la seconde comporte deux niveaux très éloignés. La moyenne seule masque cette différence.
\end{corrige}\end{blocLivre}
\begin{blocLivre}{cor-18-14}\begin{corrige}{14}{Une graduation trompeuse}\demarche{Contrôler l’échelle avant d’interpréter.}
Non : $12:8=1{,}5$ compare seulement les hauteurs au-dessus de 40. Le rapport réel est $52:48=13/12\approx1{,}083$. La différence est $\res4$ personnes. Il faut lire l’origine tronquée de l’axe.
\label{fig-cor-18-14}\figureGeo[corrige]{stat-14}
\end{corrige}\end{blocLivre}
\begin{blocLivre}{cor-18-15}\begin{corrige}{15}{Préparer une enquête}\demarche{Définir la collecte puis choisir la représentation.}
Les fréquences sont $\res{40\,\%,30\,\%,20\,\%,10\,\%}$, total 100 \%. Un diagramme en barres compare les effectifs ; un circulaire montre les parts, avec angles $144,108,72,36^\circ$. Préciser la question et le groupe interrogé.
\label{fig-cor-18-15}\figureGeo[corrige]{stat-15}
\end{corrige}\end{blocLivre}
